% ========================================================================
% Abstract
% ========================================================================
\begin{abstract}
Short-video sequential recommendation must predict a user's next interest
while simultaneously avoiding consecutive runs of highly similar items ---
a form of positional repetition that standard set-level diversity metrics
such as ILD fail to express. TRIER, a strong recent baseline for this task,
relies purely on learned ID embeddings and has no mechanism to bias the
scoring function towards locally diverse orderings. We propose PACER, a
two-stage extension of TRIER that augments item representations with
lightweight content embeddings at scoring time, then regularises the
greedy decoder with a differentiable expected-adjacent-risk objective at
training time and a small positional penalty $\lambda_c$ at inference.
Across four KuaiRec user-group partitions plus ML-1M, KuaiRand1K and
MicroLens, PACER (with $\gamma_o=0$ and $\lambda_c=0.01$) reduces
consecutive-item similarity CS@20 by 3--4$\times$ over plain TRIER on every
KuaiRec variant, drops MaxRun@20 from 12.7 to 1.0 on KuaiRec
First-Average (perfect local diversity), and improves NDCG@20 over the
same plain TRIER baseline on all four datasets. At the accuracy level,
PACER is the best performing TRIER-family model on ML-1M (HR@20=0.2736,
NDCG@20=0.1088) and the second-best on every KuaiRec variant. These results
indicate that a content-augmented scoring function paired with a
training-free positional penalty is a practical way to improve local-order
diversity without sacrificing recommendation accuracy.
\end{abstract}


% ========================================================================
% Conclusion
% ========================================================================
\begin{conclusion}
This paper investigated how short-video sequential recommendation can be
made less repetitive without compromising next-item accuracy --- a tension
that existing work tends to frame as an all-or-nothing choice between
diversity and utility. We argued that the root cause of positional
repetition is not an absence of diversity loss terms but rather the lack
of \emph{content-aware scoring}: pure ID embeddings cannot tell which
items are semantically close, so the greedy decoder defaults to
popularity-driven runs regardless of how many diversity regularisers are
added at training. To address this we introduced PACER, whose design split
into two deliberately small components so that neither introduces a new
training phase nor a new backbone. The first component augments TRIER's
scoring function with a lightweight linear projection of categorical
content embeddings, giving the decoder an item-content distance signal it
did not have before. The second component adds two complimentary
regularisation mechanisms: a differentiable expected-adjacent-risk loss
controlled by $\gamma_o$, applied only during training, and a tiny
inference-time penalty $\lambda_c$ that re-weights candidate scores by
their distance from the last chosen item. Neither mechanism is required
--- dropping $\gamma_o$ (i.e., $\gamma_o=0$) still gives the bulk of the
diversity gain, and dropping $\lambda_c$ still keeps accuracy high --- so
the user-facing trade-off stays transparent and continuous.
The most stable finding across seven datasets is that PACER cuts
consecutive-item similarity CS@20 by a factor of 3--4 over the plain
TRIER baseline on every KuaiRec variant (0.0000 on First-Average with
MaxRun@20=1.0, 0.074--0.094 on the other three). This is not a loss of
accuracy: PACER remains the best-performing TRIER-family model on ML-1M
(HR@20=0.2736, NDCG@20=0.1088) and the second-best on every KuaiRec
variant, behind the pure GRU4Rec baseline but well ahead of plain TRIER.
The Pareto sweep of \cref{sec:hyperparam} further shows that $\lambda_c$
moves the operating point along a smooth frontier: for small $\lambda_c$
values (up to about 0.015), NDCG loss stays below 0.002 while CS@20
drops by half an order of magnitude; beyond that threshold accuracy
begins to fall off more steeply, giving practitioners a clear stopping
rule. The trade-off is, however, dataset-dependent: ML-1M, with its
coarser 18-way category vocabulary and concentrated popularity
distribution, shows elevated raw repetition that benefits from a larger
$\lambda_c$ (the Pareto plot confirms that ML-1M is the outlier where
CS@20 stays above 0.50 until $\lambda_c \approx 0.02$). By contrast
KuaiRand1K and MicroLens, though extremely sparse ($\sim$1\,000 users,
2\,000--20\,000 items), still exhibit well-behaved ILD and CC scores,
indicating that the content-augmented architecture does not collapse
even when training data is scarce; the low absolute accuracy numbers on
these two datasets are consistent across all baselines and reflect a
cold-start problem rather than a PACER-specific failure mode.
Several limitations follow from the work. First, all evaluation is
offline: we measure HR, NDCG, ILD, CC, CS and MaxRun on held-out
interactions, but we cannot draw conclusions about real user
satisfaction, session retention, or long-term behavioural outcomes.
Consecutive-similarity metrics correlate with user fatigue in related
work, but a causal claim requires an online experiment. Second, the
choice of $\lambda_c$ currently relies on a small per-dataset grid sweep
(0--0.17 in 0.01 steps). A production deployment would need to automate
this, possibly by treating $\lambda_c$ as a learned hyperparameter
conditioned on user history length or item-category entropy. Third, we
have only tested PACER on sequential recommendation; whether the same
content-augmented scoring plus positional penalty transfers to
multi-task, multi-objective or graph-neural-network backbones is an open
question. Fourth, our category embeddings are derived from a simple
one-hot projection. Datasets with richer side information (tags,
transcripts, visual features) would likely benefit from stronger content
encoders, and the relative importance of $\lambda_c$ vs.\ content
quality remains to be quantified.
The natural next step, then, is an online A/B test on KuaiRand
where we compare three arms: plain TRIER, PACER with $\gamma_o=0$
but no $\lambda_c$, and PACER with $\gamma_o=0$ and a dataset-tuned
$\lambda_c$. Watching metrics of interest would be watch-through rate,
number of sessions per user, and self-reported diversity ratings collected
via an in-app prompt. A secondary analysis would correlate the observed
MaxRun@20 in logs with user drop-off rate, to test whether the offline
proxy does indeed predict real-world dissatisfaction. More broadly, the
PACER architecture suggests a principle that we suspect is under-explored:
many recommendation failures are not due to weak sequence modelling but
to scoring functions that lack the right distance signals. Whenever item
metadata is available, a small content-projection plus a train-free
positional penalty may be a cheap and general fix, not just for short-video
scenarios but for music streaming, podcast discovery, and any setting
where the user's experience depends on the local shape of the top-$K$ list.
\end{conclusion}
